% Options for packages loaded elsewhere
\PassOptionsToPackage{unicode}{hyperref}
\PassOptionsToPackage{hyphens}{url}
%
\documentclass[
]{article}
\usepackage{amsmath,amssymb}
\usepackage{iftex}
\ifPDFTeX
  \usepackage[T1]{fontenc}
  \usepackage[utf8]{inputenc}
  \usepackage{textcomp} % provide euro and other symbols
\else % if luatex or xetex
  \usepackage{unicode-math} % this also loads fontspec
  \defaultfontfeatures{Scale=MatchLowercase}
  \defaultfontfeatures[\rmfamily]{Ligatures=TeX,Scale=1}
\fi
\usepackage{lmodern}
\ifPDFTeX\else
  % xetex/luatex font selection
\fi
% Use upquote if available, for straight quotes in verbatim environments
\IfFileExists{upquote.sty}{\usepackage{upquote}}{}
\IfFileExists{microtype.sty}{% use microtype if available
  \usepackage[]{microtype}
  \UseMicrotypeSet[protrusion]{basicmath} % disable protrusion for tt fonts
}{}
\makeatletter
\@ifundefined{KOMAClassName}{% if non-KOMA class
  \IfFileExists{parskip.sty}{%
    \usepackage{parskip}
  }{% else
    \setlength{\parindent}{0pt}
    \setlength{\parskip}{6pt plus 2pt minus 1pt}}
}{% if KOMA class
  \KOMAoptions{parskip=half}}
\makeatother
\usepackage{xcolor}
\usepackage[margin=1in]{geometry}
\usepackage{color}
\usepackage{fancyvrb}
\newcommand{\VerbBar}{|}
\newcommand{\VERB}{\Verb[commandchars=\\\{\}]}
\DefineVerbatimEnvironment{Highlighting}{Verbatim}{commandchars=\\\{\}}
% Add ',fontsize=\small' for more characters per line
\usepackage{framed}
\definecolor{shadecolor}{RGB}{248,248,248}
\newenvironment{Shaded}{\begin{snugshade}}{\end{snugshade}}
\newcommand{\AlertTok}[1]{\textcolor[rgb]{0.94,0.16,0.16}{#1}}
\newcommand{\AnnotationTok}[1]{\textcolor[rgb]{0.56,0.35,0.01}{\textbf{\textit{#1}}}}
\newcommand{\AttributeTok}[1]{\textcolor[rgb]{0.13,0.29,0.53}{#1}}
\newcommand{\BaseNTok}[1]{\textcolor[rgb]{0.00,0.00,0.81}{#1}}
\newcommand{\BuiltInTok}[1]{#1}
\newcommand{\CharTok}[1]{\textcolor[rgb]{0.31,0.60,0.02}{#1}}
\newcommand{\CommentTok}[1]{\textcolor[rgb]{0.56,0.35,0.01}{\textit{#1}}}
\newcommand{\CommentVarTok}[1]{\textcolor[rgb]{0.56,0.35,0.01}{\textbf{\textit{#1}}}}
\newcommand{\ConstantTok}[1]{\textcolor[rgb]{0.56,0.35,0.01}{#1}}
\newcommand{\ControlFlowTok}[1]{\textcolor[rgb]{0.13,0.29,0.53}{\textbf{#1}}}
\newcommand{\DataTypeTok}[1]{\textcolor[rgb]{0.13,0.29,0.53}{#1}}
\newcommand{\DecValTok}[1]{\textcolor[rgb]{0.00,0.00,0.81}{#1}}
\newcommand{\DocumentationTok}[1]{\textcolor[rgb]{0.56,0.35,0.01}{\textbf{\textit{#1}}}}
\newcommand{\ErrorTok}[1]{\textcolor[rgb]{0.64,0.00,0.00}{\textbf{#1}}}
\newcommand{\ExtensionTok}[1]{#1}
\newcommand{\FloatTok}[1]{\textcolor[rgb]{0.00,0.00,0.81}{#1}}
\newcommand{\FunctionTok}[1]{\textcolor[rgb]{0.13,0.29,0.53}{\textbf{#1}}}
\newcommand{\ImportTok}[1]{#1}
\newcommand{\InformationTok}[1]{\textcolor[rgb]{0.56,0.35,0.01}{\textbf{\textit{#1}}}}
\newcommand{\KeywordTok}[1]{\textcolor[rgb]{0.13,0.29,0.53}{\textbf{#1}}}
\newcommand{\NormalTok}[1]{#1}
\newcommand{\OperatorTok}[1]{\textcolor[rgb]{0.81,0.36,0.00}{\textbf{#1}}}
\newcommand{\OtherTok}[1]{\textcolor[rgb]{0.56,0.35,0.01}{#1}}
\newcommand{\PreprocessorTok}[1]{\textcolor[rgb]{0.56,0.35,0.01}{\textit{#1}}}
\newcommand{\RegionMarkerTok}[1]{#1}
\newcommand{\SpecialCharTok}[1]{\textcolor[rgb]{0.81,0.36,0.00}{\textbf{#1}}}
\newcommand{\SpecialStringTok}[1]{\textcolor[rgb]{0.31,0.60,0.02}{#1}}
\newcommand{\StringTok}[1]{\textcolor[rgb]{0.31,0.60,0.02}{#1}}
\newcommand{\VariableTok}[1]{\textcolor[rgb]{0.00,0.00,0.00}{#1}}
\newcommand{\VerbatimStringTok}[1]{\textcolor[rgb]{0.31,0.60,0.02}{#1}}
\newcommand{\WarningTok}[1]{\textcolor[rgb]{0.56,0.35,0.01}{\textbf{\textit{#1}}}}
\usepackage{graphicx}
\makeatletter
\def\maxwidth{\ifdim\Gin@nat@width>\linewidth\linewidth\else\Gin@nat@width\fi}
\def\maxheight{\ifdim\Gin@nat@height>\textheight\textheight\else\Gin@nat@height\fi}
\makeatother
% Scale images if necessary, so that they will not overflow the page
% margins by default, and it is still possible to overwrite the defaults
% using explicit options in \includegraphics[width, height, ...]{}
\setkeys{Gin}{width=\maxwidth,height=\maxheight,keepaspectratio}
% Set default figure placement to htbp
\makeatletter
\def\fps@figure{htbp}
\makeatother
\ifLuaTeX
  \usepackage{luacolor}
  \usepackage[soul]{lua-ul}
\else
  \usepackage{soul}
\fi
\setlength{\emergencystretch}{3em} % prevent overfull lines
\providecommand{\tightlist}{%
  \setlength{\itemsep}{0pt}\setlength{\parskip}{0pt}}
\setcounter{secnumdepth}{-\maxdimen} % remove section numbering
\ifLuaTeX
  \usepackage{selnolig}  % disable illegal ligatures
\fi
\usepackage{bookmark}
\IfFileExists{xurl.sty}{\usepackage{xurl}}{} % add URL line breaks if available
\urlstyle{same}
\hypersetup{
  pdftitle={Exploratory Data Analysis},
  pdfauthor={K.D},
  hidelinks,
  pdfcreator={LaTeX via pandoc}}

\title{Exploratory Data Analysis}
\author{K.D}
\date{2026-04-21}

\begin{document}
\maketitle

\subsection{Setup}\label{setup}

\subsection{Introduction}\label{introduction}

This document highlights some preliminary EDA which uncovers interesting
patterns in the margin data.\\
Less focus has been placed on the business cycle indicators (Gross
Domestic Product and Total Hours Worked) as their general properties are
well understood.

As a reminder, \textbf{GOP margins} are defined as:

\[
\text{GOP Margin} = \frac{\text{Gross Operating Profit}}{\text{Sales}}
\]

While \textbf{GOS margins} are defined as:

\[
\text{GOS Margin} = \frac{\text{Gross Operating Surplus}}{\text{Gross Value Added}}
\]

Given that these measures are constructed from different underlying
aggregates, they display differing sensitivity to fluctuations in
business conditions, resulting in distinct margin dynamics over the
business cycle.

\includegraphics{EDA_files/figure-latex/unnamed-chunk-1-1.pdf}
\includegraphics{EDA_files/figure-latex/unnamed-chunk-1-2.pdf}

*These plots highlight the Total aggregate, which is an average of the
profit margin across all selected industries (See Appendix 1.A).

\subsection{Descriptive statistics}\label{descriptive-statistics}

\begin{Shaded}
\begin{Highlighting}[]
\FunctionTok{desc\_stats\_ts}\NormalTok{(gop\_wide, }\AttributeTok{value\_col =}\NormalTok{ Total)}
\end{Highlighting}
\end{Shaded}

\begin{verbatim}
## # A tibble: 1 x 7
##    mean    sd   min   p25 median   p75   max
##   <dbl> <dbl> <dbl> <dbl>  <dbl> <dbl> <dbl>
## 1  11.9  1.96  8.76  10.4   11.3  13.1  16.8
\end{verbatim}

\begin{Shaded}
\begin{Highlighting}[]
\FunctionTok{desc\_stats\_ts}\NormalTok{(gos\_wide, }\AttributeTok{value\_col =}\NormalTok{ Total)}
\end{Highlighting}
\end{Shaded}

\begin{verbatim}
## # A tibble: 1 x 7
##    mean    sd   min   p25 median   p75   max
##   <dbl> <dbl> <dbl> <dbl>  <dbl> <dbl> <dbl>
## 1  45.9  1.78  43.4  44.8   45.5  46.6  54.7
\end{verbatim}

\texttt{GOP} margins are lower on average (11.85\%) than \texttt{GOS}
margins (45.87\%). This may highlight more variability in operating
profits or sales over time, as they are more susceptible to changes in
economic conditions.

\subsection{Distribution Highlights}\label{distribution-highlights}

These plots track profit margins across all selected industries,
particularly highlighting the Total aggregate

\begin{verbatim}
## Warning: package 'ggrepel' was built under R version 4.5.1
\end{verbatim}

\includegraphics{EDA_files/figure-latex/pressure-1.pdf}
\includegraphics{EDA_files/figure-latex/pressure-2.pdf}

\texttt{GOP} margins across sectors are relatively more clustered than
\texttt{GOS}. This may reflect that there is smaller variation in
business profit across sectors compared to the surplus gained from
production across sectors. the \texttt{GOP} Total aggregate sits just
above the median, perhaps being skewed by the significant contributions
of the mining sector (the top-most line).

\subsection{Variance Decomposition}\label{variance-decomposition}

As we have shown, the margins are skewed by large contributing sectors.

\includegraphics{EDA_files/figure-latex/unnamed-chunk-2-1.pdf}
\includegraphics{EDA_files/figure-latex/unnamed-chunk-2-2.pdf}

\texttt{GOP} does not account for the agricultural sector, making
\texttt{Total\ less\ mining} (an aggregate subtracting the contributions
of the mining sector) the closest comparison measure.

\texttt{GOP} becomes significantly reduced and less volatile after
abstracting from the mining sector, signifying a more centralised margin
measure. The same holds true for \texttt{GOS}, as removing mining
reduces outlying spikes in the original \texttt{Total} series. Further
abstraction from agriculture only reduces the level of the series,
reflecting that the agricultural sector doesnt drive the shape of the
time-series.

This legitimises our approach of using \texttt{Total\ less\ mining} as a
true domestic margin proxy.

\subsection{Further Time-Series
Analysis}\label{further-time-series-analysis}

\begin{Shaded}
\begin{Highlighting}[]
\NormalTok{gop\_ts }\OtherTok{\textless{}{-}}\NormalTok{ gop }\SpecialCharTok{|\textgreater{}}
  \FunctionTok{mutate}\NormalTok{(}\AttributeTok{date =} \FunctionTok{yearquarter}\NormalTok{(date)) }\SpecialCharTok{|\textgreater{}}
  \FunctionTok{as\_tsibble}\NormalTok{(}\AttributeTok{key =}\NormalTok{ industry, }\AttributeTok{index =}\NormalTok{ date)}

\NormalTok{gos\_ts }\OtherTok{\textless{}{-}}\NormalTok{ gos }\SpecialCharTok{|\textgreater{}}
  \FunctionTok{mutate}\NormalTok{(}\AttributeTok{date =} \FunctionTok{yearquarter}\NormalTok{(date)) }\SpecialCharTok{|\textgreater{}}
  \FunctionTok{as\_tsibble}\NormalTok{(}\AttributeTok{key =}\NormalTok{ industry, }\AttributeTok{index =}\NormalTok{ date)}
\end{Highlighting}
\end{Shaded}

\subsubsection{Seasonal Plots}\label{seasonal-plots}

\begin{Shaded}
\begin{Highlighting}[]
\NormalTok{gop\_ts }\SpecialCharTok{|\textgreater{}} 
  \FunctionTok{filter}\NormalTok{(industry }\SpecialCharTok{==} \StringTok{"Total less mining"}\NormalTok{) }\SpecialCharTok{|\textgreater{}} 
  \FunctionTok{gg\_season}\NormalTok{(}\AttributeTok{y =}\NormalTok{ margin) }\SpecialCharTok{+} \FunctionTok{labs}\NormalTok{(}\AttributeTok{title =} \StringTok{"GOP"}\NormalTok{)}
\end{Highlighting}
\end{Shaded}

\includegraphics{EDA_files/figure-latex/unnamed-chunk-4-1.pdf}

\begin{Shaded}
\begin{Highlighting}[]
\NormalTok{gos\_ts }\SpecialCharTok{|\textgreater{}} 
  \FunctionTok{filter}\NormalTok{(industry }\SpecialCharTok{==} \StringTok{"Total less mining"}\NormalTok{) }\SpecialCharTok{|\textgreater{}} 
  \FunctionTok{gg\_season}\NormalTok{(}\AttributeTok{y =}\NormalTok{ margin)}\SpecialCharTok{+} \FunctionTok{labs}\NormalTok{(}\AttributeTok{title =} \StringTok{"GOS"}\NormalTok{)}
\end{Highlighting}
\end{Shaded}

\includegraphics{EDA_files/figure-latex/unnamed-chunk-4-2.pdf}

These seasonal plots visualise changes within a single year, overlaid
across multiple years to allow for direct comparison.

\texttt{GOP} shows a small increase across time, following the general
increase in macroeconomic variables over time. However, there is a
significant spike occurring during 2020 (highlighted by the
corresponding blue colour), the same holds true for \texttt{GOS},
highlighting that the COVID-19 period is a significant outlier for the
series.

\subsubsection{Subseries plots}\label{subseries-plots}

Subseries plots visualise the change in each quarter across the length
of the time series, and are valuable to anaylse particular outliers in a
series.

\begin{Shaded}
\begin{Highlighting}[]
\NormalTok{gop\_ts }\SpecialCharTok{|\textgreater{}} 
  \FunctionTok{filter}\NormalTok{(industry }\SpecialCharTok{==} \StringTok{"Total less mining"}\NormalTok{) }\SpecialCharTok{|\textgreater{}} 
  \FunctionTok{gg\_subseries}\NormalTok{()}
\end{Highlighting}
\end{Shaded}

\begin{verbatim}
## Plot variable not specified, automatically selected `y = gop`
\end{verbatim}

\includegraphics{EDA_files/figure-latex/unnamed-chunk-5-1.pdf}

\begin{Shaded}
\begin{Highlighting}[]
\NormalTok{gos\_ts }\SpecialCharTok{|\textgreater{}} 
  \FunctionTok{filter}\NormalTok{(industry }\SpecialCharTok{==} \StringTok{"Total less mining"}\NormalTok{) }\SpecialCharTok{|\textgreater{}} 
  \FunctionTok{gg\_subseries}\NormalTok{()}
\end{Highlighting}
\end{Shaded}

\begin{verbatim}
## Plot variable not specified, automatically selected `y = gva`
\end{verbatim}

\includegraphics{EDA_files/figure-latex/unnamed-chunk-5-2.pdf}

\begin{Shaded}
\begin{Highlighting}[]
\CommentTok{\# The change has been much more steady, with Q2 and Q3 having less of a response to COVID{-}19}
\CommentTok{\#!! shows gos is less volatile than GOP, possible connection to GOP being more responsive to business cycle (conjecture)}
\end{Highlighting}
\end{Shaded}

\texttt{GOP} has a more rigid structure in its changed across time,
reflecting higher volatility. There is a break in the series at
2015\textasciitilde, with a two significant spikes occurring during
2020, with the largest impact happening over Q3 of that year.\\
The change is much more steady for \texttt{GOS}, with Q2 and Q3 having
less of a response to COVID-19.

\subsubsection{STL Decomposition}\label{stl-decomposition}

An STL Decomposition splits a time series into three additive
components:

\[Y_t = T_t + S_t + R_t\]

Where: * \(Y_t\) (Observed Data): The raw margin percentage at time
\(t\). * \(T_t\) (Trend-Cycle): The long-term progression of the series,
capturing the structural movement of margins while smoothing out
high-frequency noise. * \(S_t\) (Seasonal): The recurring intra-year
fluctuations. \(R_t\) (Remainder): The residual component representing
stochastic ``shocks'' or noise not captured by the trend or seasonality.

\begin{Shaded}
\begin{Highlighting}[]
\NormalTok{gop\_ts }\SpecialCharTok{|\textgreater{}}
  \FunctionTok{filter}\NormalTok{(industry }\SpecialCharTok{==} \StringTok{"Total less mining"}\NormalTok{) }\SpecialCharTok{|\textgreater{}}
  \FunctionTok{model}\NormalTok{(}\FunctionTok{STL}\NormalTok{(margin)) }\SpecialCharTok{|\textgreater{}}
  \FunctionTok{components}\NormalTok{() }\SpecialCharTok{|\textgreater{}}
  \FunctionTok{autoplot}\NormalTok{() }\SpecialCharTok{+} \FunctionTok{labs}\NormalTok{(}\AttributeTok{title =} \StringTok{"GOP {-} STL Decomposition"}\NormalTok{)}
\end{Highlighting}
\end{Shaded}

\includegraphics{EDA_files/figure-latex/unnamed-chunk-6-1.pdf}

\begin{Shaded}
\begin{Highlighting}[]
\NormalTok{gos\_ts }\SpecialCharTok{|\textgreater{}}
  \FunctionTok{filter}\NormalTok{(industry }\SpecialCharTok{==} \StringTok{"Total less mining"}\NormalTok{) }\SpecialCharTok{|\textgreater{}}
  \FunctionTok{model}\NormalTok{(}\FunctionTok{STL}\NormalTok{(margin)) }\SpecialCharTok{|\textgreater{}}
  \FunctionTok{components}\NormalTok{() }\SpecialCharTok{|\textgreater{}}
  \FunctionTok{autoplot}\NormalTok{() }\SpecialCharTok{+} \FunctionTok{labs}\NormalTok{(}\AttributeTok{title =} \StringTok{"GOS {-} STL Decomposition"}\NormalTok{)}
\end{Highlighting}
\end{Shaded}

\includegraphics{EDA_files/figure-latex/unnamed-chunk-6-2.pdf}

\ul{GOP:}

\begin{itemize}
\item
  \textbf{Trend:} The series exhibits a cyclical, wavy pattern that
  remained relatively stable until a significant spike during the
  COVID-19 pandemic; it has since reverted to its pre-pandemic
  trajectory, suggesting potential for AR modeling.
\item
  \textbf{Seasonality:} The data shows severe heteroskedasticity, with
  seasonal variance expanding well before the pandemic. Because seasonal
  maxima peaked prior to COVID-19, the pandemic cannot be cited as the
  sole driver of margin expansion.
\item
  \textbf{Remainder:} Persistent patterns in the residuals indicate that
  the decomposition failed to capture all underlying variation from the
  original series.
\end{itemize}

\ul{GOS}:

\begin{itemize}
\item
  \textbf{Overall Dynamics:} The series mirrors the structural dynamics
  observed in the GOP decomposition.
\item
  \textbf{Seasonality:} Seasonal amplitude tapers toward the middle of
  the series before expanding in later periods.
\end{itemize}

Both GOP and GOS margins exhibit cyclical trends with significant
COVID-period spikes that have since normalized, alongside persistent
residual patterns and heteroskedastic seasonality that expanded well
before the pandemic.

\section{Business Cycle Indicators}\label{business-cycle-indicators}

GDP and THW follow a general exponential trend across time, mostly
following an additive (linear) process

\includegraphics{EDA_files/figure-latex/unnamed-chunk-7-1.pdf}
\includegraphics{EDA_files/figure-latex/unnamed-chunk-7-2.pdf}

\texttt{THW} appears to be more volatile, potentially due to underlying
dynamics in the labour market compared to GDP aggregates.

\subsection{Simple Transformations}\label{simple-transformations}

To make margins and the business cycle indicators comparable, we need to
transform the indicators in some way.

\paragraph{Log transformations}\label{log-transformations}

\includegraphics{EDA_files/figure-latex/unnamed-chunk-8-1.pdf}
\includegraphics{EDA_files/figure-latex/unnamed-chunk-8-2.pdf}

While this transformation does help stabilise the change in THW, it does
little to extract any cyclical variation in GDP over time.

\subsubsection{Hodrick--Prescott (HP)
Filter}\label{hodrickprescott-hp-filter}

To fully deconstruct the series into a testable cyclical component, we
need to use more advanced econometric techniques

The HP filter decomposes a time series \(y_t\) into a \textbf{trend}
component \(\tau_t\) and a \textbf{cyclical} component \(c_t\):

\[
y_t = \tau_t + c_t
\]

It does this by solving:

\[
\min_{\{\tau_t\}} \sum_{t=1}^{T} (y_t - \tau_t)^2 
+ \lambda \sum_{t=2}^{T-1} \left[ (\tau_{t+1} - \tau_t) - (\tau_t - \tau_{t-1}) \right]^2
\]

\begin{itemize}
\tightlist
\item
  First term: fit to the data\\
\item
  Second term: smoothness of the trend\\
\item
  \(\lambda\): smoothing parameter
\end{itemize}

\begin{center}\rule{0.5\linewidth}{0.5pt}\end{center}

\subsubsection{\texorpdfstring{Choice of
\(\lambda\)}{Choice of \textbackslash lambda}}\label{choice-of-lambda}

\begin{itemize}
\tightlist
\item
  Annual data: \(\lambda = 100\)\\
\item
  Quarterly data: \(\lambda = 1600\)\\
\item
  Monthly data: \(\lambda = 14400\)
\end{itemize}

\begin{center}\rule{0.5\linewidth}{0.5pt}\end{center}

\subsubsection{Interpretation}\label{interpretation}

\begin{itemize}
\tightlist
\item
  \(\tau_t\): long-run trend\\
\item
  \(c_t = y_t - \tau_t\): business cycle (deviations from trend)
\end{itemize}

\begin{Shaded}
\begin{Highlighting}[]
\FunctionTok{library}\NormalTok{(mFilter)}
\end{Highlighting}
\end{Shaded}

\begin{verbatim}
## Warning: package 'mFilter' was built under R version 4.5.3
\end{verbatim}

\begin{Shaded}
\begin{Highlighting}[]
\FunctionTok{library}\NormalTok{(tidyverse)}
\FunctionTok{library}\NormalTok{(pracma)}
\end{Highlighting}
\end{Shaded}

\begin{verbatim}
## Warning: package 'pracma' was built under R version 4.5.3
\end{verbatim}

\begin{verbatim}
## 
## Attaching package: 'pracma'
\end{verbatim}

\begin{verbatim}
## The following object is masked from 'package:purrr':
## 
##     cross
\end{verbatim}

\begin{Shaded}
\begin{Highlighting}[]
\NormalTok{hp\_gdp }\OtherTok{\textless{}{-}}\NormalTok{ mFilter}\SpecialCharTok{::}\FunctionTok{hpfilter}\NormalTok{(}\FunctionTok{log}\NormalTok{(gdp}\SpecialCharTok{$}\NormalTok{value), }\AttributeTok{freq =} \DecValTok{16000}\NormalTok{)}

\NormalTok{cycle }\OtherTok{\textless{}{-}}\NormalTok{ hp\_gdp}\SpecialCharTok{$}\NormalTok{cycle}
\NormalTok{cycle\_smooth }\OtherTok{\textless{}{-}} \FunctionTok{movavg}\NormalTok{(cycle, }\AttributeTok{n=}\DecValTok{4}\NormalTok{, }\AttributeTok{type =} \StringTok{"s"}\NormalTok{)}

\NormalTok{hp\_gdp }\OtherTok{\textless{}{-}}\NormalTok{ gdp }\SpecialCharTok{|\textgreater{}}
  \FunctionTok{mutate}\NormalTok{(}
    \AttributeTok{trend =}\NormalTok{ hp\_gdp}\SpecialCharTok{$}\NormalTok{trend,}
    \AttributeTok{cycle =}\NormalTok{ hp\_gdp}\SpecialCharTok{$}\NormalTok{cycle,}
    \AttributeTok{cycle\_smooth =}\NormalTok{ cycle\_smooth}
\NormalTok{  )}
\end{Highlighting}
\end{Shaded}

\includegraphics{EDA_files/figure-latex/unnamed-chunk-10-1.pdf}

The approximate cyclical component reflects \% deviations from the trend
of the series. In effect, it can be interpreted as a proxy for
fluctuations in business activity after holding the general increase in
wealth overtime constant. It aligns with economic intuition of peaks and
troughs,

The HP filter, along with most econometric filters, have significant
end-point problems (as shown below). This shows that it is important to
remove these values from future research as they will skew results and
do not reflect sound economic reasoning. In addition, a large trough
still occurs during the COVID-19 period. This signifies that the
COVID-19 period may need to be removed from further research, or at
least controlled for.

\includegraphics{EDA_files/figure-latex/unnamed-chunk-11-1.pdf}

\section{Appendix 1.A}\label{appendix-1.a}

\subsubsection{Industries}\label{industries}

\begin{itemize}
\item
  Agriculture
\item
  Mining
\item
  Manufacturing\\
\item
  Electricity, Gas, Water and Waste Services\\
\item
  Construction\\
\item
  Wholesale Trade\\
\item
  Retail Trade\\
\item
  Accommodation and Food Services\\
\item
  Transport, Postal and Warehousing\\
\item
  Information Media and Telecommunications\\
\item
  Financial and Insurance Services\\
\item
  Rental, Hiring and Real Estate Services\\
\item
  Professional, Scientific and Technical Services\\
\item
  Administrative and Support Services\\
\item
  Arts and Recreation Services\\
\item
  Other Services
\end{itemize}

\subsubsection{Aggregates}\label{aggregates}

\begin{itemize}
\tightlist
\item
  Total
\item
  Total less mining
\item
  Total less mining and agriculture
\end{itemize}

\end{document}
